% from ega3-1-fr.tex
\oldpage[III]{118}
\begin{proof}
La question étant locale sur $Y$, on peut supposer $Y$ noethérien, et il en
est donc de même de $X$ (I, 6.3.7). Par hypothèse, tout sous-préschéma fermé
$Z$ de $X$ dont l'espace sous-jacent est $\operatorname{Supp}(\mathcal F)$
est propre sur $Y$, autrement dit, si $j:Z\to X$ est l'injection canonique,
$f\circ j:Z\to Y$ est propre. Or, on peut supposer $Z$ tel que
$\mathcal F=j_*(\mathcal G)$, où $\mathcal G=j^*(\mathcal F)$ est un
$\mathcal O_Z$-Module cohérent (I, 9.3.5)~; comme on a
$R^qf_*(\mathcal F)=R^q(f\circ j)_*(\mathcal G)$ par
(\hyperref[III.1.3.4-fr]{1.3.4}), la conclusion résulte aussitôt de
(\hyperref[III.3.2.1-fr]{3.2.1}).
\end{proof}

\subsection{Généralisation du théorème de finitude (schémas usuels)}
\label{subsection:III.3.3-fr}

\begin{proposition}[3.3.1]
\label{III.3.3.1-fr}
Soient $Y$ un préschéma noethérien, $\mathcal S$ une
$\mathcal O_Y$-Algèbre graduée à degrés positifs, quasi-cohérente et de type
fini, $Y'=\operatorname{Proj}(\mathcal S)$ et $g:Y'\to Y$ le morphisme
structural. Soient $f:X\to Y$ un morphisme propre,
$\mathcal S'=f^*(\mathcal S)$,
$\mathcal M=\bigoplus_{k\in\mathbb Z}\mathcal M_k$ un $\mathcal S'$-Module
gradué quasi-cohérent et de type fini. Alors les
\[
R^pf_*(\mathcal M)=\bigoplus_{k\in\mathbb Z}R^pf_*(\mathcal M_k)
\]
sont des $\mathcal S$-Modules gradués de type fini pour tout $p$. Supposons
en outre que $\mathcal S$ soit engendrée par $\mathcal S_1$~; alors, pour
chaque $p\in\mathbb Z$, il existe un entier $k_p$ tel que pour tout
$k\geq k_p$ et tout $r\geq0$, on ait
\[
R^pf_*(\mathcal M_{k+r})=\mathcal S_rR^pf_*(\mathcal M_k).
\tag{3.3.1.1}\label{III.3.3.1.1-fr}
\]
\end{proposition}

\begin{proof}
La première assertion est identique à l'énoncé de
(\hyperref[III.2.4.1-fr]{2.4.1}, (i)), où on a simplement remplacé
«~morphisme projectif~» par «~morphisme propre~». Or, dans la démonstration
de (\hyperref[III.2.4.1-fr]{2.4.1}, (i)), l'hypothèse sur $f$ a été utilisée
uniquement pour montrer (avec les notations de cette démonstration) que
$R^pf'_*(\widetilde{\mathcal M})$ est un $\mathcal O_{Y'}$-Module cohérent.
Avec les hypothèses de (\hyperref[III.3.3.1-fr]{3.3.1}), $f'$ est propre
(II, 5.4.2, (iii)), donc on peut reprendre sans changement toute la
démonstration de (\hyperref[III.2.4.1-fr]{2.4.1}, (i)), grâce au théorème de
finitude (\hyperref[III.3.2.1-fr]{3.2.1}).

Quant à la seconde assertion, il suffit de remarquer qu'il y a un
recouvrement ouvert affine fini $(U_i)$ de $Y$ tel que les restrictions à
$U_i$ des deux membres de
(\hyperref[III.3.3.1.1-fr]{3.3.1.1}) soient égales pour tout
$k\geq k_{p,i}$ (II, 2.1.6, (ii))~; il suffit de prendre pour $k_p$ le plus
grand des $k_{p,i}$.
\end{proof}

\begin{corollary}[3.3.2]
\label{III.3.3.2-fr}
Soient $A$ un anneau noethérien, $\mathfrak m$ un idéal de $A$, $X$ un
$A$-schéma propre, $\mathcal F$ un $\mathcal O_X$-Module cohérent. Alors,
pour tout $p\geq0$, la somme directe
$\bigoplus_{k\geq0}H^p(X,\mathfrak m^k\mathcal F)$ est un module de type fini
sur l'anneau $S=\bigoplus_{k\geq0}\mathfrak m^k$~; en particulier, il existe
un entier $k_p\geq0$ tel que pour tout $k\geq k_p$, et tout $r\geq0$, on ait
\[
H^p(X,\mathfrak m^{k+r}\mathcal F)
=\mathfrak m^rH^p(X,\mathfrak m^k\mathcal F).
\tag{3.3.3.1}\label{III.3.3.2.1-fr}
\]
\end{corollary}

\begin{proof}
Il suffit d'appliquer (\hyperref[III.3.3.2-fr]{3.3.2}) avec
$Y=\operatorname{Spec}(A)$, $\mathcal S=\widetilde S$,
$\mathcal M_k=\mathfrak m^k\mathcal F$, compte tenu de
(\hyperref[III.1.4.11-fr]{1.4.11}).
\end{proof}

Il convient de rappeler que la structure de $S$-module de
$\bigoplus_{k\geq0}H^p(X,\mathfrak m^k\mathcal F)$ s'obtient en considérant,
pour tout $a\in\mathfrak m^r$, l'application
\[
H^p(X,\mathfrak m^k\mathcal F)
\longrightarrow H^p(X,\mathfrak m^{k+r}\mathcal F)
\]
qui provient par passage à la cohomologie de la multiplication
$\mathfrak m^k\mathcal F\to\mathfrak m^{k+r}\mathcal F$ définie par $a$
(\hyperref[III.2.4.1-fr]{2.4.1}).

